\chapter{Networking and Storage}
\label{ch:networking-storage}

\section{VNET backends}

A VNET jail owns its network stack. Blackship can attach that stack through two
native FreeBSD dataplanes:

\begin{shiptable}{L{0.25\textwidth}L{0.30\textwidth}L{0.30\textwidth}}
\toprule
Property & epair backend & netgraph backend \\
\midrule
Host switching & \command{if\_bridge} & \command{ng\_bridge} node \\
Jail interface & epair peer & \command{ng\_eiface} \\
Gateway location & bridge interface & host-side \command{ng\_eiface} \\
VLAN filtering & supported on capable kernels & not supported \\
Default & yes & no \\
\bottomrule
\end{shiptable}

Create a persisted named network rather than manually creating only a bridge:

\begin{lstlisting}[language=bash]
blackship network create frontend \
  --subnet 10.20.0.0/24 \
  --gateway 10.20.0.1 \
  --backend epair \
  --bridge blackship0
\end{lstlisting}

The record survives reboot; transient interfaces do not. \command{blackship
setup} realizes missing bridges and gateway interfaces from the record.

\section{Netgraph topology}

For a netgraph network, one host-side eiface carries the gateway and each jail
receives another eiface on a free \command{linkN} hook:

\begin{shipfigure}
\begin{tikzpicture}[
  n/.style={draw=steel,fill=codeback,rounded corners,align=center,font=\small},
  a/.style={-{Latex[length=2mm]},thick,color=signalred},node distance=0.8cm]
  \node[n] (bridge) {ng\_bridge\\ngbridge0};
  \node[n,left=of bridge] (host) {ngeth0\\host gateway\\10.20.0.1};
  \node[n,right=of bridge,yshift=0.65cm] (web) {ngeth1\\jail web\\10.20.0.10};
  \node[n,right=of bridge,yshift=-0.65cm] (db) {ngeth2\\jail db\\10.20.0.11};
  \draw[a,<->] (host) -- node[above,font=\scriptsize]{link0} (bridge);
  \draw[a,<->] (bridge) -- node[above,font=\scriptsize]{link1} (web);
  \draw[a,<->] (bridge) -- node[below,font=\scriptsize]{link2} (db);
\end{tikzpicture}
\end{shipfigure}

Outbound traffic, on either backend, requires host forwarding and a PF NAT
rule. Blackship enables forwarding while realizing a gateway but installs no
NAT: egress belongs in the host PF policy, on the interface carrying the
default route. Without it, jails reach the gateway and each other, and nothing
else:

\begin{lstlisting}
nat on $ext_if inet from 10.20.0.0/24 to any -> ($ext_if)
\end{lstlisting}

\subsection{Bridge-name collision}

Loading \command{ng\_ether} creates netgraph nodes for ethernet interfaces,
including \command{if\_bridge} devices. An epair network named
\command{blackship0} can therefore occupy that name in both namespaces. Use a
different name such as \command{ngbridge0} for an actual netgraph bridge.

\section{Address and interface ownership}

Address claims are serialized in the lease store, so concurrent creators cannot
receive the same IP. Host-side interfaces are placed in the
\command{blackship} interface group. Cleanup verifies group ownership before
destroying an interface found by name, protecting an unrelated interface after
name reuse.

The default MAC uses FreeBSD's OUI plus a stable hash of jail and bridge names.
It survives restart while differing between clones. Explicitly configure a MAC
only when an external network policy requires it.

\section{Resolver behavior for ephemeral runs}

An ephemeral \command{run} without \command{--network} inherits the host
network stack and copies the host's \file{/etc/resolv.conf}. A run attached to
a named network receives its own VNET stack. Blackship still derives its
resolver from the host, but removes loopback nameservers and interface-scoped
link-local nameservers that the VNET stack cannot reach.

If filtering leaves no usable nameserver---commonly when the host uses
\command{local\_unbound} on \command{127.0.0.1}---creation fails instead of
starting a jail with predictably broken DNS. Supply one or more reachable
addresses explicitly:

\begin{lstlisting}[language=bash]
blackship run --name probe --release 15.1-RELEASE \
  --network frontend --dns 10.20.0.53 --dns 1.1.1.1 -- fetch https://example.com
\end{lstlisting}

The repeated \command{--dns} values replace, rather than supplement, the
host-derived nameserver list. Managed manifest jails continue to use the DNS
configuration described in Chapter~\ref{ch:configuration}.

\section{PF port forwarding}

\begin{lstlisting}[language=bash]
blackship expose web --port 443 --bind-ip 192.0.2.10
blackship expose web --port 8080 --internal 80 --proto tcp
blackship ports web
blackship unexpose web
\end{lstlisting}

Blackship writes dedicated PF anchors instead of rewriting the host ruleset.
Persisted forward definitions can be reapplied with \command{setup}. Ensure the
main ruleset contains \command{rdr-anchor "blackship"} and
\command{anchor "blackship"} before publishing a service;
\command{setup --enable-pf} writes a minimal \file{/etc/pf.conf} with both when
none exists. The \command{blackship} anchor in turn references every drain
sub-anchor, so those two lines also cover draining and \setting{isolate}.

Draining uses a per-jail sub-anchor to reject new TCP SYN packets while allowing
existing states to expire. It is not a substitute for application-level
readiness removal, but closes the race between a load balancer update and jail
teardown.

\section{ZFS layout and provisioning}

With \setting{zpool = "zroot"} and \setting{dataset = "blackship"}, Blackship
organizes release and jail datasets beneath that root. Release datasets expose a
\command{@pristine} seed; jails clone from it and mount at the same
\file{data\_dir/jails/name} path used by directory-backed roots.

Advantages include:

\begin{itemize}
  \item near-instant, copy-on-write provisioning;
  \item cheap point-in-time snapshots and clones;
  \item rolling build-layer snapshots;
  \item zero-copy commit into reusable releases; and
  \item native send/receive for backup and migration.
\end{itemize}

\section{Snapshot workflow}

\begin{lstlisting}[language=bash]
blackship snapshot create web before-upgrade
blackship exec web -- pkg upgrade -y

# Rollback requires a stopped jail.
blackship down web
blackship snapshot rollback web before-upgrade --force
blackship up web

blackship clone web@before-upgrade web-test
\end{lstlisting}

Snapshot and commit operations stage configuration under \file{.blackship/}
inside the dataset so a sent snapshot carries the metadata needed to understand
its origin.

\command{commit} promotes the new release so it no longer depends on the jail.
A side effect of \command{zfs promote}: snapshots the jail took before the
commit move to the release, and the jail can no longer roll back to them.

\section{Export, import, and migration}

\begin{lstlisting}[language=bash]
# Portable archive (bsdtar): hardlinks and file flags kept, sockets
# skipped, devfs and other mounts left out.
blackship export web --output web.tar.zst

# ZFS stream.
blackship export web --output web.zfs --zfs-send

# Restore or rename.
blackship import web.tar.zst --name web-restored

# Import a foreign jail layout.
blackship import legacy.tar --from auto --name legacy

# Move a stopped jail over SSH/ZFS receive.
blackship migrate web admin@other-host
\end{lstlisting}

A tar export is written beside its destination and renamed into place, so a
failed or interrupted export never leaves a truncated archive. Foreign import
supports iocage, ezjail, plain rootfs, or automatic detection.
Review the resulting manifest and networking before first start. Migration
requires compatible ZFS and SSH authority on the destination and does not make
an incompatible FreeBSD userland/kernel combination safe. It runs as root, so
the connection uses root's keys and \file{known\_hosts}, not the invoking
user's.

\begin{shipwarning}[Know the dataset before forcing]
Forced rollback destroys newer snapshots. Forced garbage collection or manual
\command{zfs destroy -r} can remove data Blackship cannot reconstruct. List the
resolved datasets and take an export before destructive recovery.
\end{shipwarning}
